\documentclass[10pt,journal,compsoc]{IEEEtran}
\usepackage[utf8]{inputenc}
\usepackage[T1]{fontenc}
\usepackage{url}
\usepackage[hidelinks]{hyperref}
\usepackage{cite}
\usepackage{amsmath,amssymb,amsfonts}
\usepackage{algorithmic}
\usepackage{graphicx}
\usepackage{textcomp}
\usepackage{xcolor}
\usepackage{tabularx}
\usepackage{listings}
\usepackage{array}

\lstset{
  basicstyle=\small\ttfamily,
  breaklines=true,
  frame=single,
  columns=fullflexible
}

\begin{document}

\title{Arxiv Submission Guide}
\author{Bryan Alexander Buchorn \\ Independent Researcher \\ bryan.alexander@buchorn.com}
\markboth{CEREBRUM v1.2.4 Hardened Edition · March 2026}{}

\maketitle

\begin{abstract}
This module forms a critical component of the CEREBRUM framework, a community-structured graph attention engine for Knowledge Graph reasoning. It focuses on Framework Implementation Guide.
\end{abstract}

\begin{IEEEkeywords}
Knowledge Graph, Graph Attention, DSCF, CSA, Hierarchical Reasoning.
\end{IEEEkeywords}

\section{arXiv Submission Guide}

\textbf{CEREBRUM — 16 Papers Ready for Submission}

This guide covers the submission process, per-paper metadata, and checklist for submitting all sixteen CEREBRUM papers to arXiv.

---

\subsection{arXiv Categories}

All sixteen papers should be submitted to \textbf{primary category \texttt{cs.IR}} (Information Retrieval) with the following cross-list categories as appropriate:

\begin{center}
{\footnotesize
\begin{tabularx}{\columnwidth}{|>{\raggedright\arraybackslash}X|>{\raggedright\arraybackslash}X|>{\raggedright\arraybackslash}X|}
\hline
Primary & Cross-list & Use for \\ \hline
\texttt{cs.IR} & \texttt{cs.AI}, \texttt{cs.LG} & All papers (default) \\ \hline
\texttt{cs.AI} & \texttt{cs.IR}, \texttt{cs.NE} & Papers 3, 4, 7 (neuro\-morphic / biological analog) \\ \hline
\texttt{cs.LG} & \texttt{cs.IR}, \texttt{cs.AI} & Papers 6, 15 (learning algorithms) \\ \hline
\texttt{cs.DC} & \texttt{cs.IR}, \texttt{cs.AI} & Papers 5, 13 (distributed / streaming) \\ \hline
\end{tabularx}
}
\end{center}

---

\subsection{Submission Order (Recommended)}

Submit found\-ational papers first so cross-references resolve correctly.

\begin{center}
{\footnotesize
\begin{tabularx}{\columnwidth}{|>{\raggedright\arraybackslash}X|>{\raggedright\arraybackslash}X|>{\raggedright\arraybackslash}X|>{\raggedright\arraybackslash}X|}
\hline
Order & Paper & File & Category \\ \hline
1 & DSCF/TSC Community Fusion & \texttt{PAPER\_001\_DSCF\_TSC.md} & \texttt{cs.IR} + \texttt{cs.AI} \\ \hline
2 & Community-Structured Attention & \texttt{PAPER\_002\_CSA.md} & \texttt{cs.IR} + \texttt{cs.LG} \\ \hline
3 & Bridge Twin Nodes & \texttt{PAPER\_003\_BRIDGE\_TWINS.md} & \texttt{cs.AI} + \texttt{cs.NE} \\ \hline
4 & STDP Causal Edge Inference & \texttt{PAPER\_004\_STDP\_CAUSAL.md} & \texttt{cs.AI} + \texttt{cs.NE} \\ \hline
5 & Holographic Federated Index & \texttt{PAPER\_005\_HOLOGRAPHIC\_INDEXING.md} & \texttt{cs.DC} + \texttt{cs.IR} \\ \hline
6 & Bayesian Beam Search & \texttt{PAPER\_006\_BAYESIAN\_BEAM.md} & \texttt{cs.LG} + \texttt{cs.IR} \\ \hline
7 & REM Cycle Maintenance & \texttt{PAPER\_007\_REM\_CYCLE.md} & \texttt{cs.AI} + \texttt{cs.IR} \\ \hline
8 & Cross-Modal Signal Alignment & \texttt{PAPER\_008\_SIGNAL\_ENCODER.md} & \texttt{cs.IR} + \texttt{cs.LG} \\ \hline
9 & THALAMUS Ingestion Pipeline & \texttt{PAPER\_009\_THALAMUS.md} & \texttt{cs.IR} + \texttt{cs.AI} \\ \hline
10 & Inference Validation & \texttt{PAPER\_010\_INFERENCE\_VALIDATION.md} & \texttt{cs.IR} + \texttt{cs.AI} \\ \hline
11 & Contra\-diction Materi\-alization & \texttt{PAPER\_011\_CONTRADICTION.md} & \texttt{cs.IR} + \texttt{cs.AI} \\ \hline
12 & Glass-Box Reasoning Studio & \texttt{PAPER\_012\_REASONING\_STUDIO.md} & \texttt{cs.HC} + \texttt{cs.IR} \\ \hline
13 & Streaming Engine & \texttt{PAPER\_013\_STREAMING\_ENGINE.md} & \texttt{cs.DC} + \texttt{cs.IR} \\ \hline
14 & Metaco\-gnitive Verifi\-cation & \texttt{PAPER\_014\_INSIGHT\_ENGINE.md} & \texttt{cs.AI} + \texttt{cs.IR} \\ \hline
15 & Algorithmic Depth & \texttt{PAPER\_015\_ALGORITHMIC\_DEPTH.md} & \texttt{cs.LG} + \texttt{cs.IR} \\ \hline
16 & Production Hardening & \texttt{PAPER\_016\_PRODUCTION\_HARDENING.md} & \texttt{cs.SE} + \texttt{cs.IR} \\ \hline
\end{tabularx}
}
\end{center}

---

\subsection{Per-Paper Metadata Template}

Each submission requires the following metadata fields on arXiv:

\begin{lstlisting}
Title:        [From paper header]
Authors:      Bryan Alexander Buchorn; Claude Sonnet 4.6 (Research Collaborator, Anthropic)
Abstract:     [From paper Abstract section — max 1,920 characters]
Comments:     16 pages. Part of the CEREBRUM framework series.
              Code: https://github.com/[repo]
MSC-class:    68T30 (Knowledge representation)
ACM-class:    I.2.4; H.3.3
License:      CC BY 4.0 (arXiv default for non-commercial research)
\end{lstlisting}

---

\subsection{Formatting Requir\-ements}

arXiv accepts \texttt{.tex} (preferred) and \texttt{.pdf}. Our papers are in Markdown; conversion is required.

\subsubsection{Markdown $\to$ LaTeX conversion}

\begin{lstlisting}
pip install pandoc
pandoc docs/arxiv/PAPER_001_DSCF_TSC.md -o arxiv_submission/paper_001.tex \
    --template arxiv_template.tex \
    --bibliography refs.bib
\end{lstlisting}

A LaTeX template (\texttt{arxiv\_template.tex}) following arXiv's \texttt{article} document class should be prepared once and reused for all sixteen papers.

\subsubsection{Required LaTeX elements}
\begin{itemize}
\item \texttt{\textbackslash usepackage{amsmath}} — for displayed equations
\item \texttt{\textbackslash usepackage{booktabs}} — for benchmark tables
\item \texttt{\textbackslash usepackage{hyperref}} — for cross-references and URLs
\item All figures as \texttt{.pdf} or \texttt{.png} at 300 DPI minimum

\end{itemize}
\subsubsection{Figure requi\-rements}
Each paper should include at minimum:
\begin{enumerate}
\item An archi\-tecture diagram (flowchart of the algorithm)
\item A results table (benchmark comparison)

\end{enumerate}
Figures should be prepared in a vector format (PDF) for crisp rendering.

---

\subsection{Abstract Length Check}

arXiv abstracts are capped at 1,920 characters. Check each abstract:

\begin{lstlisting}
for paper in docs/arxiv/PAPER_*.md; do
    abstract=$(awk '/^### Abstract/{found=1; next} found && /^###/{exit} found{print}' "$paper")
    count=$(echo "$abstract" | wc -c)
    echo "$paper: $count chars"
done
\end{lstlisting}

---

\subsection{Pre-Submission Checklist (per paper)}

\begin{itemize}
\item [ ] Abstract $\le$ 1,920 characters
\item [ ] All equations render correctly in LaTeX
\item [ ] All tables have \texttt{\textbackslash toprule/\textbackslash midrule/\textbackslash bottomrule} (booktabs style)
\item [ ] All citations in \texttt{refs.bib} with correct BibTeX keys
\item [ ] Author affil\-iations: "Independent Researcher" + "Anthropic"
\item [ ] Prior art section explicitly diffe\-rentiates from closest prior work
\item [ ] Benchmark numbers match \texttt{release/BENCHMARKS.md}
\item [ ] GitHub/code link included in Comments field
\item [ ] License selected: CC BY 4.0
\item [ ] Cross-list categories correct for paper topic

\end{itemize}
---

\subsection{Citation Cross-Referencing}

Papers cite each other. Use consistent arXiv IDs once papers 1–4 are submitted. Placeholder format until IDs are assigned:

\begin{lstlisting}
@article{buchorn2026dscf,
  title   = {Dual-Signal Community Fusion for Knowledge Graph Attention},
  author  = {Buchorn, Bryan Alexander and {Claude Sonnet 4.6}},
  journal = {arXiv preprint arXiv:2026.XXXXX},
  year    = {2026}
}
\end{lstlisting}

---

\subsection{Compre\-hensive Paper (Parallax\_White\_Paper\_arXiv.md)}

The compr\-ehensive umbrella paper (\texttt{docs/Parallax\_White\_Paper\_arXiv.md}) covers all subsystems in one document and should be submitted to \texttt{cs.IR} as the primary reference after all sixteen individual papers are live. Title it:

\textgreater  \textbf{CEREBRUM: Community-Structured Graph Attention for Zero-Shot Multi-Hop Knowledge Graph Reasoning}

This paper should cite all sixteen subsystem papers by their arXiv IDs.

---
\textbf{Copyright © 2026 Bryan Alexander Buchorn. All Rights Reserved.}


\bibliographystyle{IEEEtran}
\bibliography{../../docs/latex/references}

\end{document}
